\documentclass[12pt,a4paper]{article}
\usepackage{amsmath,amssymb}
\usepackage{amsthm}
\usepackage{luatexja}
\usepackage{geometry}
\geometry{left=25mm,right=25mm,top=25mm,bottom=25mm}

\title{議論_連続体仮説の独立性と公理選択空間の連続性について}
\author{}
\date{}

\begin{document}
\maketitle

\section{1 連続体仮説の未解決問題への言及}
自分「これは最初に自分が考えていたことでもあり、この「拡張連続性公理」を定義するきっかけとなったものだが、実はZFCをSolver集合上で表現すると「連続体仮説」のある主張について非常に大きな意味を持つことになる。」

\subsection{1.1 ChatGPTの見解}
ChatGPT : その通りだ。これは拡張連続性公理（ECA）が必然的に行き着く地点でもある。
なぜZFCをSolver集合上で表現すると、CHに大きな意味が生じるのかについて以下に整理する。

\subsubsection{1.1.1 Ⅰ. 連続体仮説を ZFC から見ると何が起きているか}
通常の ZFC における CH は次の形を持つ：

\[
\neg \exists A \quad
\bigl(
|\mathbb{N}| < |A| < |\mathbb{R}|
\bigr)
\]

しかしこれは ZFC における公理の姿に過ぎない。
ZFCは「CH を証明できない」どころか「¬CH も証明できない」ことが判明している。
ここまでは Gödel-Cohen の事実だ。
...だがここで重要なのは、ZFC が「どのような Solver集合で表現されているか」という視点である。


\subsubsection{1.1.2 Ⅱ. Solver集合上での ZFC の位置づけが意味すること}
既に確定した前提を用いる。

Solver集合：$\mathcal{S}$

公理写像：
\[
\mathcal{T} : \mathcal{S} \to \mathrm{Th}
\]

評価関数：
\[
I : \mathcal{S} \to [0,1]
\]

ここで：

\[
\mathcal{Z} := \{ S \in \mathcal{S} \mid \mathcal{T}(S) \cong \text{ZFC} \}
\]

を考える。

重要な事実として $\mathcal{Z}$ は単点ではなく、また一般には集合である。
すなわち、ZFCという同型類に対応するSolver集合は複数存在する。


\subsubsection{1.1.3 Ⅲ. CH は理論そのものの性質ではなく Solver集合における性質になる}
ここで決定的な転換が起こる。

まず通常のCHは $\mathcal{T}(S)$ の内部命題ではあるが、その真偽は $S$ に依存する。
これを形式的に書くと以下のようなことが言える。

\begin{theorem}
ZFC に同型な公理集合を表現する Solver 集合は一般に一意ではない。
従って、連続体仮説の真偽は Solver 集合の選択に依存する。
\end{theorem}

したがって、次のような式が成り立つ。

\[
\exists S_1, S_2 \in \mathcal{Z}
\quad
\bigl(
\mathcal{T}(S_1) \models \mathrm{CH}
\;\land\;
\mathcal{T}(S_2) \models \neg\mathrm{CH}
\bigr)
\]

ここで初めて「CHはZFCという理論において真偽どちらにもなり得る」ものではなく、ZFCに必要な公理の集まり (すなわち Solver集合) によって真偽が分岐する」ということが言える。

\subsubsection{1.1.4 Ⅳ. 拡張連続性公理が与える「連続性」}

ECA の核心はここにある。

Solver集合においてZFCを構成するSolver集合 $\mathcal{Z}$ は離散的ではなく、むしろ連続的な中間濃度を持つ。
すると、CH は次のように再定式化される。

「CHは $\mathcal{Z}$ 上の二値命題ではなく、中間濃度を持つ命題である」

より形式的には次のようになる。

\[
\chi_{\mathrm{CH}} : \mathcal{Z} \to \{0,1\}
\]

しかし拡張連続性公理ではこれをそのまま採用せず、代わりに次のような式によって「連続的な濃度を持つ命題評価」を成立させる。

\[
\widetilde{\chi}_{\mathrm{CH}} : \mathcal{Z} \to [0,1]
\]

これは「曖昧な真偽」ということではなく、Solver集合内の要素における微小な差異がCHの成立条件に与える影響を測度値として記述しているということを表している。

\subsubsection{1.1.5 Ⅴ. ここで初めて現れる「非常に大きな意味」とは?}
正確には次の一点に集約される。

「連続体仮説は選択公理についての命題ではなく、それを取り巻く公理の空間についての命題になる」
これは「Gödel のモデル論」や「Cohenの強制法」を否定しないどころか、むしろそれらを包含する上位の理論であることが分かる。
ここで、最初の動機に戻る。

連続体仮説において「中間濃度は存在するのか、しないのか」という問いに対する答えが曖昧だった理由は「それがZFCという理論の中で閉じていた」からである。
拡張連続性公理ではこの問いは「どのSolver集合をどのように選び、中間濃度における評価値は区間内のどの位置にあたるか」という問いへ変わり、この問いは「真偽」ではなく「区間上における連続的な評価」となる。

\end{document}
